\chapter[1920 --- Krogh]{1920: August Krogh}
\label{ch:1920_augustkrogh}

\section*{Main Contribution}

\textit{Discovered the mechanism by which capillaries regulate blood flow to tissues, explaining how the body matches oxygen delivery to metabolic demand.}

\section{Official Citation}

for his discovery of the mechanism of regulation of respiration in animals

\section{Background and Historical Context}

Schack August Steen Krogh was born on November 15, 1874, in Grenaa, a small fishing town on the east coast of Jutland, Denmark. The son of a shipbuilder, Krogh grew up surrounded by the sea, and his early fascination with marine biology and the physiology of aquatic organisms would ultimately lead him to the frontiers of respiratory physiology. After completing his secondary education at the Grenaa Gymnasium, Krogh enrolled at the University of Copenhagen in 1893, where he studied zoology and physiology. He was initially drawn to marine biology and spent several summers studying the respiration of aquatic invertebrates at the marine laboratory at Helsing{\o}r, but his interests soon shifted to the broader questions of how organisms acquire and transport oxygen---questions that would define his scientific career.

At the University of Copenhagen, Krogh came under the influence of the distinguished physiologist Christian Bohr, who was then conducting pioneering research on the binding of oxygen to hemoglobin. Bohr, together with his colleagues Peter Astrup and Kaj Roholm, had demonstrated that the oxygen-binding curve of hemoglobin was sigmoidal (S-shaped), a property that was essential for the efficient loading of oxygen in the lungs and its unloading in the tissues. The ``Bohr effect''---the dependence of hemoglobin's oxygen affinity on pH and carbon dioxide concentration---was one of a major discoveries in respiratory physiology, and it provided the framework within which Krogh would develop his own insights into gas exchange.

The early twentieth century was a period of intense interest in respiratory physiology. Scientists understood the basic anatomy of the lungs and the general principles of gas exchange---that oxygen diffused from the inspired air into the blood, and that carbon dioxide diffused from the blood into the expired air---but the mechanisms by which oxygen moved from the alveoli into the blood and from the blood into the tissues, and by which these processes were regulated, remained poorly characterized. The prevailing view, influenced by the work of the German physiologist Hermann von Helmholtz and the Danish physiologist Christian Bohr, held that respiration was a passive process governed by simple diffusion gradients. Krogh would challenge this simplistic view and demonstrate that the respiratory system was far more dynamic and precisely regulated than anyone had imagined.

The timing of Krogh's Nobel Prize was unusual. He received it in 1920 for work that had been published a decade earlier, and the award was actually given for the year 1919 (there were no prizes awarded during 1915--1918 due to World War I). The delay reflected both the disruption caused by the war and the Nobel Committee's recognition that Krogh's work, while initially received with some skepticism, had been confirmed and extended by numerous investigators in the intervening years. The award recognized not only Krogh's specific discoveries about capillary regulation but also his broader contributions to the methodology and philosophy of physiology.

\section{The Discovery and Contribution}

Krogh's most celebrated achievement was his discovery of the mechanism by which capillary blood flow is regulated to match the metabolic demands of tissues---a phenomenon now known as ``capillary recruitment'' or the ``Krogh cylinder'' model of oxygen delivery. To appreciate the significance of this discovery, it is necessary to understand the state of knowledge about capillary function at the time.

At the turn of the twentieth century, the capillaries---the smallest blood vessels, with diameters of only 5--10 micrometers---were generally regarded as passive conduits for blood flow. The prevailing view, held by most physiologists, was that capillaries were always open and that blood flow through them was determined solely by the pressure gradient between arteries and veins. This view was based on the assumption that the capillary walls, being only a single cell thick, had no capacity for active contraction or relaxation. Krogh's experiments would overturn this assumption and reveal the capillary as a dynamic, regulated structure.

Krogh's approach was characteristically innovative. He developed new techniques for observing and measuring capillary blood flow in living animals, using the tail of the eel, the tongue of the frog, and the skeletal muscles of various mammals as experimental models. By carefully counting the number of open capillaries in resting and exercising muscles, Krogh made a remarkable observation: in resting muscle, only a fraction of the available capillaries were open to blood flow, but upon exercise, many previously closed capillaries opened up, dramatically increasing the surface area available for gas exchange. This observation---that the number of perfused capillaries could be regulated---was a paradigm shift.

Krogh demonstrated that the increase in capillary surface area during exercise was not simply a passive consequence of increased blood flow, but rather an active process mediated by the contraction and relaxation of the smooth muscle cells surrounding the capillaries (the precapillary sphincters). He showed that this capillary recruitment was regulated by local metabolic factors, particularly the concentration of carbon dioxide and oxygen in the surrounding tissue. When tissue oxygen levels fell or carbon dioxide levels rose (as occurred during muscular exercise), the precapillary sphincters relaxed, opening previously closed capillaries and increasing blood flow to the active tissue. Conversely, when tissue oxygen levels rose and carbon dioxide levels fell (as occurred during rest), the sphincters contracted, reducing capillary perfusion and conserving blood flow for other tissues.

Krogh formalized these observations in a mathematical model that described oxygen diffusion from a capillary into the surrounding tissue. The ``Krogh cylinder'' model treated each capillary as the center of a cylindrical region of tissue, and calculated the oxygen tension as a function of distance from the capillary and time. The model predicted that the oxygen supply to tissue depended critically on the density of perfused capillaries, and that capillary recruitment could dramatically extend the range of tissue oxygenation. The Krogh cylinder model remains a foundational concept in tissue oxygenation and microvascular physiology, and it continues to be used and refined by researchers today---for example, in the prediction of tumor hypoxia and in the design of oxygen delivery strategies for critical care.

Beyond capillary recruitment, Krogh made several other important contributions to respiratory physiology. He was among the first to apply Fick's principle of diffusion to biological systems, providing quantitative descriptions of gas exchange in the lungs and tissues. He developed the ``indirect Fick method'' for measuring cardiac output---a technique based on the analysis of oxygen consumption and arteriovenous oxygen difference that remained in clinical use for decades. He also studied the respiratory exchange of insects and the ventilation mechanisms of fish gills, demonstrating the universality of respiratory principles across the animal kingdom.

Krogh's experimental methodology was as important as his specific discoveries. He was a master of the carefully designed experiment, and he insisted on rigorous quantitative analysis of all physiological data. His textbook, \emph{The Respiratory Exchange of Animals and Man} (1916), was a landmark in the field, combining elegant experimental techniques with sophisticated mathematical analysis. The book established the standard for how respiratory physiology should be studied---a standard that Krogh maintained throughout his career. He also developed new instruments for measuring gas exchange, including the Krogh spirometer and the Krogh gas analyzer, which became standard tools in physiology laboratories worldwide.

Krogh was also a pioneer in the application of the principle of ''symmorphosis''---the idea that biological structures are designed to meet but not exceed functional demands. Although this principle was formally articulated by Krogh's student Edward Taylor and others in the 1960s, its intellectual origins lie in Krogh's work on the relationship between capillary density and tissue oxygen consumption. Krogh showed that the capillary supply of different tissues was precisely matched to their metabolic demands, a finding that suggested a remarkable degree of optimization in the design of the microcirculation.

\section{Impact on Modern Medicine}

Krogh's discoveries have had a significant impact on modern medicine, particularly in the fields of critical care, sports medicine, and vascular biology. The concept of capillary recruitment---that the microvasculature can dynamically adjust blood flow to match tissue metabolic demands---is now a fundamental principle of cardiovascular physiology. It underlies our understanding of how the body maintains tissue oxygenation during exercise, how it responds to hemorrhage and shock, and how it adapts to chronic conditions such as heart failure and chronic obstruct pulmonary disease.

The Krogh cylinder model of tissue oxygenation remains the starting point for virtually all theoretical analyses of oxygen delivery to tissues. Modern refinements of the model, incorporating factors such as heterogeneous capillary spacing, red blood cell transit time, and metabolic heterogeneity, have extended Krogh's original framework to address clinically relevant problems. These include the prediction of tissue hypoxia in tumors (which has implications for cancer therapy and radiation treatment planning), the assessment of oxygen delivery in critically ill patients, and the optimization of oxygen supplementation in neonates. The concept of ``critical capillary spacing''---the maximum distance between perfused capillaries at which tissue oxygenation remains adequate---is a direct extension of Krogh's model and is used to predict the onset of tissue ischemia in clinical settings.

In clinical practice, Krogh's insights into capillary regulation inform the management of circulatory shock. The recognition that capillary perfusion can be impaired even when gross blood flow is adequate (a condition known as microcirculatory dysfunction) has led to the development of goal-directed hemodynamic therapies that aim to optimize not just cardiac output and blood pressure, but also microvascular flow. Techniques such as sublingual video-microscopy (using devices like the MicroScan), which allows direct visualization of capillary blood flow in critically ill patients, are direct clinical applications of the principles that Krogh elucidated a century ago. The ability to monitor microcirculatory function at the bedside has transformed the management of sepsis and other forms of circulatory shock.

Krogh's work also has important implications for understanding and treating diabetes mellitus. Diabetic microangiopathy---the pathological thickening of capillary basement membranes that is a hallmark of diabetes---impairs capillary function and contributes to the devastating vascular complications of the disease, including retinopathy (which causes blindness), nephropathy (which causes kidney failure), and neuropathy (which causes nerve damage). The recognition that capillary dysfunction is central to diabetic complications has driven research into therapies that target the microcirculation, including drugs that improve capillary perfusion and reduce endothelial dysfunction. The development of angiotensin-converting enzyme (ACE) inhibitors and angiotensin receptor blockers (ARBs), which are used to slow the progression of diabetic nephropathy, is based in part on the understanding that microvascular function is a critical determinant of organ damage in diabetes.

In sports medicine, Krogh's principle of capillary recruitment is central to understanding the physiological adaptations to training. Regular exercise increases capillary density in skeletal muscle, enhancing oxygen delivery and exercise capacity. This adaptation, which can be quantified using modern imaging techniques such as contrast-enhanced ultrasound and near-infrared spectroscopy, is one of the key mechanisms by which physical training improves cardiovascular health and reduces mortality. The concept of ``capillary density'' as a determinant of aerobic fitness is a direct application of Krogh's work.

Krogh's contributions to respiratory physiology have also influenced the development of respiratory support technologies. Mechanical ventilators, which are used to support patients with respiratory failure, are designed to optimize the matching of ventilation and perfusion in the lungs---a principle that is directly related to Krogh's work on the regulation of gas exchange. The development of supplemental oxygen therapy, pulse oximetry, and transcutaneous oximetry (which measures tissue oxygenation non-invasively) all build upon the foundational understanding of tissue oxygenation that Krogh established.

\section{Legacy}

August Krogh's legacy is one of methodological innovation and conceptual insight. He demonstrated that the smallest blood vessels in the body were not passive tubes but dynamic, regulated structures that played an active role in tissue homeostasis. His insistence on rigorous quantitative analysis set standards for physiology that remain exemplary, and his development of new experimental techniques and instruments advanced the capabilities of the entire field.

After receiving the Nobel Prize, Krogh continued to make important contributions to physiology and ecology. He studied the relationship between oxygen supply and the distribution of animal species, pioneering the field of ecological physiology (now called physiological ecology). He demonstrated that the distribution of animal species in nature was determined by their physiological adaptations to the oxygen availability in their environment, and he used this principle to predict the effects of environmental changes on animal populations. This work anticipated the modern fields of ecological physiology and conservation biology, which study the relationships between organisms and their environments.

Krogh was also a passionate advocate for open science. He argued that scientific knowledge should be freely available to all, and he was instrumental in establishing the Scandinavian Physiological Society and promoting scientific collaboration across national borders. His vision of science as a collaborative enterprise, rather than a competitive race, remains an important ideal in an era of increasing specialization and competition.

The capillary, once regarded as the humblest of blood vessels, is now recognized as a critical determinant of tissue health and disease. August Krogh, through his meticulous experiments and elegant theoretical analyses, revealed the capillary to be a sophisticated and dynamic organ in its own right. His work continues to illuminate the remarkable mechanisms by which the body delivers oxygen to its tissues, and his legacy endures in every clinical assessment of microvascular function, every mathematical model of tissue oxygenation, and every study of the remarkable plasticity of the cardiovascular system. Krogh died in Copenhagen on September 13, 1949, but his influence on physiology and medicine continues to resonate through the work of the countless investigators who have built upon his foundational discoveries.


\section{Modern Developments}

Krogh's capillary physiology has informed modern vascular medicine. Understanding of capillary recruitment and angiogenesis has led to therapies for peripheral artery disease and heart failure. Anti-VEGF therapies (bevacizumab, ranibizumab) for age-related macular degeneration target the same vascular growth mechanisms. The COVID-19 pandemic highlighted the importance of capillary function, as endothelial damage and microvascular thrombosis are major complications of severe infection.
